\section{Efficient implementation: what the derivation permits}\label{sec:implementation}

The exact reference algorithm is fully specified above. This section describes
the implementable baseline, and then the three refinements that measurement
showed to matter. The baseline is not a straw man: it is the configuration the
solver still offers as its verifiable reference, and every refinement below is a
separate option so that its effect can be measured in isolation.

\subsection{A linear-memory reference organization}
Store the original directed graph in contiguous endpoint and adjacency arrays,
with incoming and outgoing access. Store the forest separately through parent,
parent-edge, root and child information. Maintain the nonbasic vertex/slack
sets and a mapping between variable IDs and their positions in basic lists.
Traverse trees iteratively, so a long chain cannot exhaust the call stack.

A postorder traversal computes $p(D_e)$ and $w(D_e)$ for every forest edge.
Component totals provide root reduced costs. Only $n-1$ nonbasic variables
need pricing. After selecting an entry, mark its subtree or tree, form
\eqref{eq:direction}, and scan all arcs for decreasing basic slacks.
Rebuild the forest metadata after the exchange. This deliberately simple
version provides a reliable reference before local updates are introduced.

\subsection{Specialized operations in one iteration}
The forest supplies specialized operations for entering-variable selection,
the pivot column, the leaving-variable test, and the basis exchange.
Table~\ref{tab:operation-costs} separates these costs for the full-rebuild
reference organization. In particular, the $O(n)$ pricing bound is independent
of graph density; the complete direction and ratio test also inspect arcs.

\begin{table}[htbp]
\centering\small
\renewcommand{\arraystretch}{1.15}
\begin{tabular}{p{.29\linewidth}p{.49\linewidth}l}
\toprule
Operation & Forest calculation & Cost\\
\midrule
Subtree aggregates & Postorder sums of profits and weights & $O(n)$\\
One reduced cost & Signed aggregate difference, sums available & $O(1)$\\
Full pricing and entry & Scan the $n-1$ nonbasic candidates & $O(n)$\\
Pivot direction on vertices & Mark a tree/subtree; normalize its motion & $O(n)$\\
Pivot direction on slacks & Head-minus-tail vertex directions & $O(m)$\\
Ratio test and exit & Scan decreasing basic vertex/slack variables & $O(n+m)$\\
Chosen pivot row, if needed & Evaluate its $n-1$ nonbasic coefficients from forest membership and aggregates & $O(n)$\\
Basis exchange and rebuild & Exchange edge/root membership; rebuild metadata & $O(n+m)$\\
\bottomrule
\end{tabular}
\caption{Arithmetic costs with the current forest available and reusable
arrays. Constant-time coordinate evaluation assumes the required sums and
membership information have been computed.}\label{tab:operation-costs}
\end{table}

\begin{proposition}[Arithmetic cost of the full-rebuild variant]
\label{prop:cost}
Using adjacency lists and reusable arrays, one full rebuild, complete pricing,
direction computation, full ratio test and exchange can be performed in
$O(n+m)$ arithmetic operations and $O(n+m)$ storage. An oracle performing
K pivots costs $O((K+1)(n+m))$ arithmetic operations.
\end{proposition}
\begin{proof}
Forest construction and traversals touch $O(n+m)$ stored entries. Subtree
sums and pricing touch $O(n)$ entries, since a forest has fewer than n edges.
Constructing vertex directions is linear in n, and evaluating all slack
directions and blockers is linear in m. Updating membership and rebuilding
the adjacency information is linear as well. Arrays store only vertices, arcs,
membership, aggregates and temporary vectors.
\end{proof}
This is an arithmetic-operation bound, not a bit-complexity bound for growing
rational numbers. Nor does it bound K polynomially. The exact checker
deliberately builds dense matrices and therefore does \emph{not} implement
the memory bound in \cref{prop:cost}.
For residual calls indexed by r, the analogous total is
\[
 O\!\left(\sum_r(K_r+1)(n_r+m_r)\right),
\]
when residual graphs are rebuilt linearly at each call. The number of calls
is at most the original n; no parametric-flow complexity bound is being
transferred to the \FS{}.

\subsection{Variants to implement and tune}
\begin{table}[htbp]
\centering\small
\begin{tabular}{p{.22\linewidth}p{.32\linewidth}p{.35\linewidth}}
\toprule
Parameter family & Candidate variants & Invariant to preserve\\
\midrule
Pricing & Full scan; blocked scan; candidate list & A full check before claiming no improving entry\\
Entering rule & Bland; first or best improving & A documented anti-cycling policy for each alternative\\
Forest update & Full; local; adaptive rebuild & Same mathematical basis and direction as a full reconstruction\\
Rebuild trigger & Pivot interval; changed-tree fraction; residual threshold & Recover accuracy without silently changing the problem\\
Initial basis & Sink singleton; closure of a promising seed; best undirected component & Primal feasibility: the positive component must be closed in the residual\\
Ratio test & Full scan; restricted to boundary-crossing arcs; stop at the first zero step & Same step, and the same leaving variable whenever Bland's rule is active\\
Residual start & Fresh sink basis; validated warm start & Primal feasibility on the new residual graph\\
Vertex order & Input; topological; seeded & Identical certified objective and canonical decomposition\\
\bottomrule
\end{tabular}
\caption{Tuning choices of the implementation. Every row is an explicit
parameter, so an ablation can switch off one refinement at a time.}
\end{table}
Reducing the number of visited vertices or tested candidates is useful only if
it reduces total runtime at the required accuracy. The profiler separates
pricing, directions, ratio tests, exchange, rebuilds and canonicalization. The
measured ranking of these choices is reported in the experimental report, not
here; the one result worth repeating is that the entering rule barely moves the
outcome while the treatment of degeneracy moves it by factors.

\paragraph{A caution about termination.}
First-improving or best-improving pricing is not automatically Bland's rule.
Switching to a finite rule permanently after a documented trigger, or using
a separately justified lexicographic implementation, are possible safeguards.
An informal fallback that repeatedly switches back to an unproved rule does
not inherit the finite-pivot theorem without additional reasoning. The
reasoning is supplied in \cref{prop:episodes} below, because a permanent switch
turns out to be expensive.

\subsection{Degeneracy is the dominant cost}\label{sec:degeneracy-cost}

In a basis the vertex values are constant on the positive component and zero on
every other tree. A step is therefore blocked at length zero whenever the
blocking variable is one of the many already at zero: a vertex of the entering
subtree whose direction component is negative, or a precedence slack between two
vertices that are both at zero. Degeneracy is not incidental here; it is what the
precedence structure produces.

The effect is large enough to reorganize the implementation around it. On
applied instances more than $99\%$ of the pivots of an oracle have zero step; on
a family with no arcs at all, where no slack can block, the share is zero. Three
consequences follow, and each was measured.

\paragraph{The anti-cycling fallback must not become the operating rule.}
A permanent switch to Bland's rule after a run of degenerate pivots engages
immediately and never releases, because the run never ends: every entering rule
then collapses onto Bland within the first oracle. Releasing the fallback at the
first pivot with a positive step, and placing the trigger far away, keeps the
improving rule without giving up finiteness.

\begin{proposition}[Finiteness with a released fallback]\label{prop:episodes}
Suppose the solver uses an arbitrary entering rule, switches to Bland's rule
after a fixed number of consecutive zero-step pivots, and returns to the
original rule at the first pivot with a positive step. Then the oracle performs
finitely many pivots.
\end{proposition}
\begin{proof}
Call an \emph{episode} a maximal run of iterations during which Bland's rule is
active. Within an episode every pivot has zero step, so the basic solution never
changes, and the iterations of the episode are exactly the iterations of Bland's
rule applied to the linear program from the basis at which the episode began;
by the finite-pivot theorem for Bland's rule the episode cannot cycle and must
end. An episode ends only with a pivot of positive step, which strictly
increases the objective of the normalized subproblem. Since there are finitely
many bases, hence finitely many objective values, only finitely many episodes
can occur. Between episodes only finitely many pivots are possible for the same
reason.
\end{proof}

\paragraph{The ratio test need not scan every arc.} By
\eqref{eq:direction} the direction is $+\sigma$ on the entering subtree,
$-\kappa$ on the positive component and zero elsewhere. An arc whose two
endpoints lie outside both sets therefore has zero direction difference and
cannot block. It suffices to scan the arcs incident to the entering subtree,
together with the arcs that cross the boundary of the positive component; and
since the positive component is closed, such an arc has its tail outside it, so
scanning the outgoing arcs of the vertices outside the positive component completes
the cover.

\paragraph{A zero step is the minimum, so the scan may stop at the first one.}
No step is negative, so the first blocking variable attaining zero already
determines the step. The shortcut must be suspended while Bland's rule is
active, since that rule requires the smallest index among ties. This is the
technique the group's 2023 prototype implemented as an explicit feasibility
check during the forest traversal, and it is the single largest saving in the
ratio test: on one oracle of a mining instance with $1060$ vertices the arcs
examined fall from $1\,110\,883$ to $184$, at an unchanged number of pivots.

\subsection{The dual side}\label{sec:dual}

Primal degeneracy and dual degeneracy are different defects, and on this problem
they are wildly asymmetric. Primal degeneracy is a per-pivot phenomenon: more
than $99\%$ of the pivots of an oracle have zero step. The only source of dual
degeneracy is a tie between ratios, which is a per-block phenomenon: on a
generated instance of two thousand vertices, ties affect $123$ of $267$ blocks. Even
if every tie produced one degenerate dual pivot, that is a hundred against half a
million. The dual is therefore the natural method here, which is also what the
classical guidance says: use the dual when the primal is heavily degenerate.

Two structural facts make the dual unusually simple on this problem.

\begin{proposition}[Primal infeasibility is combinatorial]\label{prop:dual-shape}
In any basis of the normalised subproblem, the vertex variables satisfy $y\ge 0$.
Every primal infeasibility is a precedence slack $s_{ij}=y_j-y_i<0$, that is an
arc leaving the positive component $P$; and the basis is primal feasible exactly
when $P$ is closed.
\end{proposition}
\begin{proof}
The basic solution takes the value $1/w(P)$ on $P$ and zero on every other tree,
so $y\ge0$ always. For an arc $(i,j)$, $s_{ij}=y_j-y_i$ is negative only when
$y_i>y_j$, which forces $i\in P$ and $j\notin P$. Conversely if $P$ is closed no
arc leaves it, all such slacks vanish or are positive, and the basis is feasible.
\end{proof}

Choosing the leaving variable is therefore a search for an unsatisfied
prerequisite, with no numerical scan, and the stopping rule is that $P$ is
closed. The starting basis is dual feasible by construction: make every vertex its
own tree and let $P$ be the vertex of maximum ratio; the reduced cost of any other
vertex $k$ is then $w_k(p_k/w_k-\rho)\le0$, and there are no forest edges yet. It
costs $O(n)$, and it is essentially forced: a singleton start is dual feasible
only for the maximum-ratio vertex.

The row of the leaving slack is as local as the primal direction. Entering a
nonbasic $q$ by one unit moves every vertex by
$\sigma_q\ind[u\in S_q]-\kappa_q\ind[u\in P]$, so the row of the slack of
$(t,h)$ is nonzero only for the root of each endpoint's tree, for the forest
edges on the two paths to their roots, and through the rank-one term $\kappa_q$
when exactly one endpoint lies in $P$. It costs $O(n+\mathrm{depth})$.

\paragraph{What the measurement says.} The dual removes the degeneracy it was
predicted to remove --- the share of zero-step pivots falls from above $99\%$ to
between $0\%$ and $21\%$ --- but the gain in time is between $1.1$ and $2.7$,
not the two orders of magnitude a naive extrapolation would suggest. The reason
is worth recording: the primal performs about half a million pivots of which only
some fifteen hundred move the objective, while the dual performs three hundred
thousand of which most do. It has traded many null steps for many tiny ones.
Choosing the leaving arc by structural criteria rather than by smallest index
--- most attractive head, largest tree, most negative reduced cost --- was tried
and does not help, because every violated slack has the same value $-1/w(P)$ and
there is nothing for a most-infeasible rule to discriminate.

\paragraph{What remains.} With those three changes the ratio test is no longer
the bottleneck. Pricing and forest rebuilding are: on a layered DAG with dense
connections between levels they account for roughly a third and a half of the
solver's time respectively, because both are $\Theta$ of the touched part of the
graph at every pivot while the pivot itself changes very little. Incremental
pricing, updated only on the components a degenerate pivot touches, and updates
confined to the modified paths, are the natural next steps; neither is
implemented here.

\subsection{Numerical checks and output contracts}
Floating-point implementations must use declared absolute and relative
tolerances, inspect primal and dual residuals in the original units, and
reconstruct values when incremental updates drift. Near-equal ratios require
particular care: approximate merging must not be mistaken for an exact tie.
On integer/rational review cases, exact comparisons provide a useful check.

The intended outputs are: ratio and connected support for an oracle;
canonical blocks and breakpoints for a path computation; a feasible x and
validated bound/certificate for a capacity solve. Time, iteration or numerical
limits must produce distinct states rather than an ``optimal'' label.
Optional internal checks can vary with configuration, but final verification
remains mandatory for results used in a scientific comparison.

Once the positive path is known, finding a capacity's split block costs
$O(\log k)$ with cumulative weights. Returning only the objective or a compact
block description can be correspondingly cheap. Materializing all n vertex
coordinates still costs $O(n)$. The many-capacity experiment must include the
initial path computation and distinguish these output requirements.
